FaceCon is a research and development project that explores how much of a person's facial shape can be predicted from their genetic data. Estimating a face from genetic material is relevant to forensic investigation, anthropology, and medical research, and recent progress in artificial intelligence has made the idea newly practical. At the same time, the most visible existing system was built on private data that other researchers cannot check or reuse. Independent experts have also questioned whether its results genuinely reflect individual genetics rather than general characteristics such as age, gender, and ancestry. FaceCon responds to this situation by building an open and carefully evaluated alternative.

The purpose of FaceCon is to provide a transparent and repeatable system that predicts the three dimensional shape of a face from genetic data, using openly documented research datasets rather than private ones. Every step of the process, from preparing the genetic data to measuring the results, is recorded so that others can repeat and verify the work. This openness is the core value of the project, because it gives the research community a trustworthy reference point in a field where strong claims have so far been difficult to confirm.

The proposed solution follows a deliberately cautious strategy. FaceCon begins with very simple prediction methods and moves to more advanced ones only if the simpler methods prove insufficient. Importantly, it always compares its predictions against a basic estimate built only from a person's age, gender, and ancestry. This comparison is what allows the project to separate genuine genetic information from general population patterns and to report honestly how much the genetic data really adds. A fully acceptable and scientifically useful outcome of the project is the finding that genetics adds little beyond this basic estimate under the available data.

Alongside the three dimensional shape of the face, the system also provides a set of visible colour traits, namely eye colour, hair colour, and skin colour, by integrating an established forensic DNA phenotyping method that reads the same genetic data. FaceCon does not train its own method for these colours; it reuses the validated external one and shows its results on the predicted face. The expected outcomes of the project are an open and reusable prediction system, a clear and honest account of facial-shape prediction accuracy and the integrated appearance predictions, and a simple demonstration tool that shows the predicted face next to the real one. The project's limitations are acknowledged from the outset. The genetic and facial data are sensitive and controlled-access, usable records are few, and the data mainly represent certain population groups. For these reasons, FaceCon is positioned as a careful research baseline rather than a finished product, and it does not claim to reconstruct any individual's exact face.

At the time of this deliverable, approval to use the real research data had not yet been granted. The team therefore built the complete system and tested it on a simulated dataset designed to resemble the real study, in which the true answer is known in advance. On this simulated data the system correctly finds genetic information when it is present and correctly reports none when it is absent, which shows that the method can be trusted. The tests were then repeated using real human genetic data from a public reference collection, with the faces still simulated, so that the genetic side of the problem carries the full difficulty of real data. The system behaved in the same way, with the expected reduction in how much it could find. Across all of these tests a simple statistical model predicted better than the more complex neural networks, even when the genetic effects were deliberately made more complicated than a simple model can represent, which suggests that at this data size the way the genetic data is presented matters more than the choice of model. These tests confirm that the system works as designed. They do not yet show how well real faces can be predicted, which is the main work planned for the second phase of the project.
